\documentclass[twocolumn]{iopjournal}

\usepackage{amsmath,amssymb,amsfonts,amsthm}

\fancyhfoffset[L]{0mm}
\setlength{\marginparwidth}{0pt}
\setlength{\marginparsep}{0pt}

\renewcommand{\articletype}[1]{%
  {\vspace*{-8mm}\noindent \Large \sf Preprint}\par
  \vspace*{6mm}\noindent{\scriptsize \sf{\bfseries \MakeUppercase{#1}}}}

\fancyhead[L]{}
\fancyhead[R]{}

\hypersetup{pdftitle={The SST Covariant Completion: y From the Acceleration of the Time Leg}, pdfauthor={Stephen Nelson}, pdfsubject={Gravity Field Dynamics (GFD)}, pdfkeywords={scalar-tensor transition, covariant action, tetrad, time leg, acceleration, khronometric, preferred foliation, MOND, proof}}

\begin{document}

\articletype{Supplementary Material}

\title{The SST Covariant Completion: $y$ From the Acceleration of the Time Leg}

\author{Stephen Nelson}

\affil{Field Dynamics Research}

\ead{stephen@fielddynamics.org}

\begin{abstract}
The quasi-static proof of the SST~\cite{NelsonPqs} left three points
open: the covariant definition of $y$, which the archive writes as
$|\nabla\Phi|^2/a_0^2$; $F$ for timelike gradients, $y<0$; and the
metric field equations, with the status of $\Phi$'s stress--energy.
We settle all three with one definition. Let $u = e_0$ be the time
leg of the tetrad and $a^\mu = u^\nu\nabla_\nu u^\mu$ its
acceleration, and set $y = c^4 a_\mu a^\mu/a_0^2$. This $y$ is a
spacetime scalar. Since $a\cdot u = 0$, the acceleration is spacelike
or zero, so $y \ge 0$ for every configuration and $F$ needs no
extension. In a static metric $-N^2c^2dt^2 + h_{ij}dx^idx^j$ with
$N = e^{\Phi/c^2}$, $c^2 a_i = \partial_i\Phi$ exactly, so $y$ reduces
to the archive's $|\nabla\Phi|^2/a_0^2$, and $\Phi$ is the lapse
potential of the metric, not a separate field. The SST, with the
reading of Lemma~1 of~\cite{NelsonPqs}, is then the action
$(1/16\pi G)\sqrt{-g}\,[R - 2f(a)]$ with $f = a_0^2[F(y) - y]$. This
is the preferred-foliation form of Blanchet and
Marsat~\cite{BM2011,BM2012}, and it gives $\mu_{\rm total} = 1 +
f'(a)/2a = x/(x+1)$ exactly. Its static weak field gives $\Psi =
\Phi$ and $\nabla\cdot[\mu\nabla\Phi] = 4\pi G\rho$; the metric
equations carry the TetB share as phantom density; the time-leg
stress--energy is post-Newtonian; and the homogeneous background is
exactly general relativity. We also state a limit the covariant form
makes precise: $f$ does not tend to a constant at large $a$, so on
each mass's own time leg the action leaves a TetB tail $\to a_0$. The
archive excludes that tail in the solar system and applies $\mu$ to
the galaxy's net field. The covariant statement of that rule is that
$a$ is the acceleration of the coarse-grained time leg; deriving that
from an action is open.
\end{abstract}

\keywords{scalar-tensor transition, covariant action, tetrad, time leg, acceleration, khronometric, preferred foliation, MOND, proof}

% ----------------------------------------------------------------------

\section*{Setting}

The SST is the action~\cite{NelsonL3,NelsonL3supp,NelsonII}
\begin{equation}
\begin{aligned}
S = \!\int\! d^4x\sqrt{-g}\Big[&\frac{R}{16\pi G} - \frac{a_0^2}{8\pi G}F(y)\\
&- \tfrac14 e^{-2\Phi/c^2}F_{\mu\nu}F^{\mu\nu} + \mathcal L_{\rm m}\Big],
\end{aligned}
\label{eq:sst}
\end{equation}
with $F(y) = y - 2\sqrt y + 2\ln(1+\sqrt y)$ and
$\mu(x) = x/(x+1)$, $x = \sqrt y$. The quasi-static
proof~\cite{NelsonPqs} reads the action as the framework does
(Lemma~1 there): in the static weak field the curvature term supplies
the Newtonian part $y$ of $F$, and the scalar term supplies
$F(y) - y$. It derives the field equation, $\gamma = 1$ and the
galactic results from that reading. Its Scope listed three open
points, and this note answers them.

\textbf{The time leg.} The tetrad's leg $e_0$ is unit timelike,
$g(e_0,e_0) = -1$. We take it normal to GFD's time foliation (in the
static case, the slices of constant Killing time), so $u = e_0$ is
hypersurface-orthogonal. Its acceleration is
$a^\mu = u^\nu\nabla_\nu u^\mu$.

\textbf{Definition.} \(y \equiv c^4\, a_\mu a^\mu / a_0^2.\)

Each identity below is checked exactly by the accompanying script, 16
checks in all, using sympy for every algebraic
step~\cite{NelsonScript}.

\textbf{Contributions and prior work.} An action of the form
$(1/16\pi G)\sqrt{-g}[R - 2f(a)]$, built from the acceleration of a
unit timelike field normal to a preferred foliation, is Blanchet and
Marsat's modified gravity with a preferred time
foliation~\cite{BM2011,BM2012}. Its MOND phenomenology, and the
conditions on $f$, are theirs. Related constructions are Sanders'
Ho\v{r}ava-type MOND~\cite{Sanders2011}, generalized Einstein-aether
MOND~\cite{ZFS2007} and the khronometric theory~\cite{BPS2011}. What
this note adds:
\begin{itemize}
\item the identification of the SST's $y$ with the acceleration of the
tetrad's time leg;
\item the SST's own $f$, $f = a_0^2[F(y) - y]$, which gives the
archive's $\mu = x/(x+1)$ exactly;
\item the identification of $\Phi$ with the lapse;
\item the strong-field limit of this $f$, and what it means for the
archive's net-field rule;
\item machine checks of every step.
\end{itemize}

% ----------------------------------------------------------------------

\section*{Theorem 1 (covariance and positivity)}

\emph{$y$ is a spacetime scalar, and $y \ge 0$ for every configuration.
So $F(y)$ is defined everywhere, and no branch for timelike gradients
is needed.}

\begin{proof}
$y$ is the contraction of the vector $a^\mu$ with itself, so it is a
scalar. From $u\cdot u = -1$, $u_\mu\nabla_\nu u^\mu = 0$, so
$a\cdot u = 0$: $a$ is orthogonal to a timelike vector and is
therefore spacelike or zero. Hence $a_\mu a^\mu \ge 0$. The script
checks $g(u,u) = -1$ and $a\cdot u = 0$ identically for a general
time-dependent lapse and conformal factor, and $a\cdot a \ge 0$ at 20
random points of a time-dependent metric.
\end{proof}

So open point~2 does not arise. An argument $\nabla_\mu\Phi\nabla^\mu\Phi$
of a scalar $\Phi$ can be negative; the time leg's acceleration
cannot.

% ----------------------------------------------------------------------

\section*{Theorem 2 (the static reduction)}

\emph{For a static metric $-N^2c^2dt^2 + h_{ij}dx^idx^j$, the time
leg's acceleration is $a_0 = 0$, $a_i = \partial_i\ln N$ exactly. With
$N = e^{\Phi/c^2}$,
\begin{equation}
c^2 a_i = \partial_i\Phi,\qquad
y = \frac{h^{ij}\partial_i\Phi\,\partial_j\Phi}{a_0^2},
\label{eq:static}
\end{equation}
and in the weak field, $h_{ij} = (1 - 2\Psi/c^2)\delta_{ij}$,
$y = |\nabla\Phi|^2/a_0^2 + O(c^{-2})$.}

\begin{proof}
Direct computation of $u^\nu\nabla_\nu u^\mu$ with
$u = (Nc)^{-1}\partial_t$ (checked for general $N$ and a conformal
$h$). Substituting $N = e^{\Phi/c^2}$ gives~(\ref{eq:static}), and the
limit $c\to\infty$ gives the weak-field form (checked).
\end{proof}

So the archive's $y$ is the static form of a covariant scalar. The
field strength is the acceleration of the frame's time leg: what an
observer at rest in the foliation measures on an accelerometer. In
the static case $\Phi \equiv c^2\ln N$ is the lapse potential of the
metric. It is not a field added to the metric, which is the framework's
reading~(i) in~\cite{NelsonPqs}, now stated exactly.

% ----------------------------------------------------------------------

\section*{Theorem 3 (the covariant SST)}

\emph{With $y$ as defined, the SST read as in Lemma~1
of~\cite{NelsonPqs} is the action
\begin{equation}
\begin{aligned}
S &= \int d^4x\,\frac{\sqrt{-g}}{16\pi G}\big[R - 2f(a)\big] + S_{\rm EM} + S_{\rm m},\\
f &= a_0^2\big[F(y) - y\big],
\end{aligned}
\label{eq:cov}
\end{equation}
in units $c = 1$, $y = a^2/a_0^2$. Then:
\begin{enumerate}
\item Blanchet and Marsat's $\chi = f'(a)/2a$ is
$\chi = \mu - 1$, so $\mu_{\rm total} = 1 + \chi = \mu(a/a_0) = x/(x+1)$
exactly;
\item the curvature term supplies the $y$ of $F$, and the time-leg
term supplies $F - y$: $-2f/16\pi G = -(a_0^2/8\pi G)[F(y) - y]$;
\item in the static weak field, with minimal coupling,
$\nabla^2(\Phi - \Psi) = 0$, so $\Psi = \Phi$ ($\gamma = 1$), and
\begin{equation}
\nabla\cdot\big[\mu(|\nabla\Phi|/a_0)\nabla\Phi\big] = 4\pi G\rho .
\label{eq:field}
\end{equation}
\end{enumerate}}

\begin{proof}
(1) Differentiate $f(a) = a_0^2[F(a^2/a_0^2) - a^2/a_0^2]$; using
$F'(y) = \mu(\sqrt y)$ gives $f'/2a = \mu - 1$ (checked exactly).
(2) is the definition. (3) By Theorem~2 the static weak-field
Lagrangian of~(\ref{eq:cov}) is Lemma~1 of~\cite{NelsonPqs},
\[
\frac{(\nabla\Psi)^2 - 2\nabla\Phi\cdot\nabla\Psi}{8\pi G}
- \frac{a_0^2}{8\pi G}\big[F(y) - y\big] - \rho\Phi .
\]
Its Euler--Lagrange equations are checked: the $\Psi$-equation is
$\nabla^2(\Phi - \Psi) = 0$, and the $\Phi$-equation on $\Psi = \Phi$
is~(\ref{eq:field}).
\end{proof}

Equation~(\ref{eq:cov}) is the reading of Lemma~1 written as a
covariant action. Read with $F$ in place of $F - y$, the $y$ would be
counted twice, and the field equation would be
$\nabla\cdot[(1+\mu)\nabla\Phi] = 4\pi G\rho$~\cite{NelsonPqs}.

% ----------------------------------------------------------------------

\section*{Theorem 4 (the metric equations and $\Phi$'s stress--energy)}

\emph{Varying~(\ref{eq:cov}) gives Einstein's equation with the
time-leg stress tensor as an extra source, plus the foliation's own
equation~\cite{BM2011}. There is no separate scalar stress tensor,
because $\Phi$ is the lapse. In the static weak field:
\begin{enumerate}
\item the $00$-equation is $\nabla^2\Phi = 4\pi G(\rho + \rho_{\rm ph})$,
with $\rho_{\rm ph} = \nabla\cdot[(1-\mu)\nabla\Phi]/4\pi G$, the
TetB share as phantom density (for a deep point mass,
$\rho_{\rm ph} = \sqrt{GMa_0}/4\pi G r^2$);
\item the time-leg energy density is the gravitational field energy,
$\sim g^2/8\pi G$, since $F - y = -y + \tfrac23 y^{3/2} + O(y^2)$ in
the deep field. As a mass density it is $v^2/2c^2$ times
$\rho_{\rm ph}$, where $v^4 = GMa_0$: about $2\times10^{-7}$ at
$200$~km/s.
\end{enumerate}}

\begin{proof}
(1) is~(\ref{eq:field}) rearranged, and the deep point-mass value is
checked for $g = \sqrt{GMa_0}/r$. (2) The series of $F - y$ and the
ratio are checked exactly.
\end{proof}

This answers open point~3. ``$\Phi$ is the metric's Newtonian
potential'' is Theorem~2 ($\Phi = c^2\ln N$) together with
Theorem~3(3) ($\Psi = \Phi$). Its energy enters the metric at
post-Newtonian order, as the energy of any gravitational field does.

% ----------------------------------------------------------------------

\section*{Theorem 5 (the cosmological background)}

\emph{In a homogeneous, isotropic background the comoving time leg is
geodesic, $a = 0$, so $y = 0$ and $F(0) = 0$. The background is
exactly general relativity, and the SST acts on perturbations.}

\begin{proof}
For $-dt^2 + a(t)^2d\mathbf x^2$ with $u = \partial_t$, every
component of $u^\nu\nabla_\nu u^\mu$ vanishes (checked), and
$F(0) = 0$.
\end{proof}

% ----------------------------------------------------------------------

\section*{Theorem 6 (the two limits of $f$)}

\emph{In the deep field, $f = -a^2 + 2a^3/3a_0 + O(a^4)$:
Blanchet and Marsat's MOND condition with $\Lambda_0 = 0$. In the
strong field, $f = -2a_0a + 2a_0^2\ln(a/a_0) + O(1)$, which does not
tend to a constant, their condition for recovering general
relativity. So the TetB share $(1-\mu)g$ tends to $a_0$ as
$x\to\infty$.}

\begin{proof}
Both series are checked exactly, as is
$\lim_{x\to\infty} a_0\mu(x) = a_0$.
\end{proof}

\textbf{What this means for the solar system.} Evaluated on each
mass's own time leg, the action leaves a TetB tail
$(1-\mu)g \to a_0$ around every star. At Saturn, $x\approx 5\times10^5$
and the tail is $a_0$ to two parts in $10^6$ (checked). This is the
per-mass tail that the archive records as excluded: a constant pull
$\approx a_0$ at every planet, ruled out by perihelia, and a
Sun-centred quadrupole, ruled out by Cassini (\cite{NelsonColes},
\S5.4; \cite{NelsonBimetric}, Part~3). The archive therefore applies
$\mu$ to the galaxy's net field, never to a star's own field. The
covariant form makes that rule precise: $a$ is the acceleration of
the coarse-grained time leg, the frame of the galaxy's net field, and
bound subsystems gravitate internally through $R$ alone. In a solar
system that net field is uniform to $\sim10^{-12}$ and is removed by
free fall~\cite{NelsonColes}. Equation~(\ref{eq:cov}) with $a$ taken
from each mass's own time leg does not implement this. An action that
does, for example by defining the coarse-graining covariantly, is the
open step.

% ----------------------------------------------------------------------

\section*{Scope}

Proved here:
\begin{itemize}
\item the covariant $y$ and its positivity;
\item its exact static reduction, and $\Phi$ as the lapse;
\item the covariant form~(\ref{eq:cov}) of the SST under Lemma~1, with
$\mu = x/(x+1)$ exactly;
\item $\gamma = 1$ and the field equation from it;
\item the structure and size of the metric equations' extra source;
\item the general-relativistic background;
\item the two limits of $f$.
\end{itemize}

Not proved here:
\begin{enumerate}
\item \emph{The net-field mechanism} (Theorem~6).
\item \emph{Preferred-frame effects.} A preferred foliation breaks
local Lorentz invariance. The PPN parameters $\alpha_1, \alpha_2$ and
the khronon's dynamics and stability~\cite{BPS2011} are not computed
for this $f$. In Blanchet and Marsat's model, general relativity is
recovered well above $a_0$ because $f$ tends to a constant there
(the cosmological constant); this $f$ does not, so these must be
computed with the net-field rule in place.
\item \emph{Signal speeds.} The action has no $(\nabla u)^2$ shear or
expansion terms, so the curvature term carries the tensor modes. The
scalar-sector propagation speed found in the archive assumed $\Phi$ as
a covariant scalar and needs redoing with $\Phi$ as the lapse.
\item \emph{The fourth term.} With $\Phi = c^2\ln N$, the EM coupling
$e^{-2\Phi/c^2}$ is defined relative to the foliation's normalized
time; its coefficient $\alpha$ remains open, as in~\cite{NelsonPqs}.
\item \emph{A general time leg.} If $e_0$ is not hypersurface-orthogonal,
the theory is of generalized Einstein-aether type~\cite{ZFS2007}, with
a vorticity sector not treated here.
\end{enumerate}

% ----------------------------------------------------------------------

\begin{thebibliography}{99}

\bibitem{NelsonPqs}
S.~Nelson,
\textit{The SST Covariant Completion in Its Quasi-Static Limit} (2026),
\texttt{proofs/sst-quasi-static-limit/}.

\bibitem{NelsonL3}
S.~Nelson,
\textit{Gravity Field Dynamics: Field Closure of the Tetrad Frame},
Letter~3 (2026).

\bibitem{NelsonL3supp}
S.~Nelson,
\textit{Covariant completion}, supplementary note to Letter~3 (2026).

\bibitem{NelsonII}
S.~Nelson,
\textit{A Covariant Theory Without Dark Matter}, Paper~II (2026),
doi:10.5281/zenodo.18360354.

\bibitem{NelsonColes}
S.~Nelson,
\textit{Gravity Field Dynamics --- Coles Notes} (2026), \S5.4.

\bibitem{NelsonBimetric}
S.~Nelson,
\textit{Bimetric projection} (2026),
\texttt{research/bimetric-projection/}, Part~3.

\bibitem{NelsonScript}
S.~Nelson,
\texttt{verify\_sst\_covariant\_time\_leg.py} (2026),
\texttt{proofs/sst-covariant-time-leg/}.

\bibitem{BM2011}
L.~Blanchet and S.~Marsat,
Modified gravity approach based on a preferred time foliation,
arXiv:1107.5264 (2011).

\bibitem{BM2012}
L.~Blanchet and S.~Marsat,
Relativistic MOND theory based on the Khronon scalar field,
arXiv:1205.0400 (2012).

\bibitem{Sanders2011}
R.~H.~Sanders,
Hiding Lorentz invariance violation with MOND,
arXiv:1105.3910 (2011).

\bibitem{ZFS2007}
T.~G.~Zlosnik, P.~G.~Ferreira and G.~D.~Starkman,
Modifying gravity with the aether: an alternative to dark matter,
\textit{Phys.\ Rev.\ D} \textbf{75} (2007) 044017.

\bibitem{BPS2011}
D.~Blas, O.~Pujol\`as and S.~Sibiryakov,
Models of non-relativistic quantum gravity: the good, the bad and the
healthy, \textit{JHEP} \textbf{04} (2011) 018.

\end{thebibliography}

\end{document}
